\section{Software Desenvolvido e Resultados}
\label{sec:resultados}

% Capturas empilhadas (largura total da página).
% Enquanto os PNG não existirem, usa example-image-16x9 (pacote mwe).
% Substituir salvando em figuras/:
%   tela-inventario.png      — dashboard / inventário
%   tela-mapa.png            — mapa geral de Contagem
%   tela-mapa-detalhe.png    — mapa com popup ou filtro aplicado
\newcommand{\imagemtela}[1]{%
  \IfFileExists{figuras/#1.png}%
    {\includegraphics[width=\linewidth,height=7.2cm,keepaspectratio]{figuras/#1.png}}%
    {\includegraphics[width=\linewidth,height=7.2cm,keepaspectratio]{example-image-16x9}}}

\subsection{A aplicação web}

A aplicação tem duas telas, o Inventário e o Mapa, acessadas por uma barra de navegação.
Em ambas, uma faixa superior informa a data da fotografia das bases, o total de estabelecimentos e quantos têm coordenada.
O Inventário e o Mapa compartilham os mesmos filtros, de modo que um filtro aplicado em uma tela vale também na outra (RNF03).

No Inventário, o visitante busca por nome, CNPJ ou bairro e filtra por subclasse da CNAE, porte, condição de MEI e disponibilidade de coordenada.
Um painel lateral mostra, em barras, o total por subclasse e por porte, e cada barra funciona como filtro.
Os resultados aparecem em uma tabela com nome, CNPJ, CNAE, porte, bairro e indicação de coordenada, limitada aos 200 primeiros registros.
A Figura \ref{fig:tela-inventario} apresenta essa tela.

\begin{figure}[htbp]
    \centering
    \caption{Tela Inventário da aplicação}
    \label{fig:tela-inventario}
    \vspace{4pt}
    \imagemtela{tela-inventario}
    \par\vspace{4pt}\textbf{\footnotesize Fonte: Elaborado pelo autor.}
\end{figure}

No Mapa, cada estabelecimento com coordenada é um marcador.
A Figura \ref{fig:tela-mapa} mostra a distribuição dos pontos em Contagem, e a Figura \ref{fig:tela-mapa-detalhe} mostra o detalhe ao selecionar um estabelecimento: o clique abre uma janela com nome, CNAE, CNPJ e o aviso de que a posição é aproximada.

\begin{figure}[htbp]
    \centering
    \caption{Tela Mapa da aplicação --- visão geral}
    \label{fig:tela-mapa}
    \vspace{4pt}
    \imagemtela{tela-mapa}
    \par\vspace{4pt}\textbf{\footnotesize Fonte: Elaborado pelo autor.}
\end{figure}

\begin{figure}[htbp]
    \centering
    \caption{Tela Mapa da aplicação --- detalhe do estabelecimento}
    \label{fig:tela-mapa-detalhe}
    \vspace{4pt}
    \imagemtela{tela-mapa-detalhe}
    \par\vspace{4pt}\textbf{\footnotesize Fonte: Elaborado pelo autor.}
\end{figure}

\subsection{Base de dados obtida e perfil dos estabelecimentos}

A integração das bases resultou em 781 estabelecimentos de TI ativos em Contagem na fotografia de 13/10/2026.
A Tabela \ref{tab:perfil} mostra a distribuição por subclasse da CNAE e por porte.
Três subclasses (suporte técnico, consultoria e programas sob encomenda) reúnem 646 estabelecimentos, ou 82,7\% do total, e as microempresas são 86,6\%, o que indica um setor formado sobretudo por pequenos prestadores de serviço.

\begin{table}[htbp]
    \centering
    \caption{Estabelecimentos de TI ativos em Contagem por subclasse da CNAE e por porte}
    \label{tab:perfil}
    \small
    \begin{tabular}{llrr}
        \hline
        \textbf{Critério} & \textbf{Categoria} & \textbf{Qtd.} & \textbf{\%} \\
        \hline
        Subclasse & 6209-1/00 Suporte técnico e outros serviços & 276 & 35,3 \\
        & 6204-0/00 Consultoria em TI & 191 & 24,5 \\
        & 6201-5/01 Programas sob encomenda & 179 & 22,9 \\
        & 6202-3/00 Programas customizáveis & 82 & 10,5 \\
        & 6203-1/00 Programas não customizáveis & 37 & 4,7 \\
        & 6201-5/02 \textit{Web design} & 16 & 2,0 \\
        \hline
        Porte & Microempresa & 676 & 86,6 \\
        & Empresa de pequeno porte & 69 & 8,8 \\
        & Demais & 36 & 4,6 \\
        \hline
    \end{tabular}
    \par\vspace{3pt}\textbf{\footnotesize Fonte: Elaborado pelo autor, com dados da Receita Federal (13/10/2026).}
\end{table}

Quanto à concentração espacial, o bairro com mais estabelecimentos é o Eldorado, com 105 (13,4\%), e os bairros que levam Eldorado no nome somam 173 (22,2\%).
Em seguida vêm Riacho das Pedras e Jardim Riacho das Pedras, com 56 (7,2\%) em conjunto, entre 125 bairros distintos.

Quanto ao MEI, nenhum dos 781 estabelecimentos foi marcado como tal: 630 empresas constavam no arquivo do Simples Nacional, e nenhuma tinha a opção pelo MEI.
Como o campo não foi validado em outra fonte, essa pergunta ficou respondida apenas em parte.

\subsection{Localização dos estabelecimentos}

Dos 781 estabelecimentos, 641 (82,1\%) receberam coordenada e 140 (17,9\%) ficaram sem ela porque o CEP não consta no CNEFE.
Todas as coordenadas atribuídas vieram da junção apenas por CEP: nenhum estabelecimento foi ligado por endereço completo.
A causa provável, não verificada, é que o logradouro do CNPJ traz o tipo (``RUA PROFESSOR ALVES HORTA''), enquanto o campo lido do CNEFE traz só o nome, e as chaves de junção não coincidem.

Quanto ao nível de geocodificação do CNEFE, 605 pontos (94,4\%) são coordenadas originais do Censo 2022 e 36 (5,6\%) são coordenadas modificadas, o que não significa que sejam a posição do estabelecimento, já que o ponto é o de um endereço do mesmo CEP.

Quanto à sobreposição, os 641 estabelecimentos ocupam apenas 360 coordenadas distintas: 103 coordenadas são compartilhadas por mais de um estabelecimento, que somam 384 (59,9\% dos georreferenciados), e a maior reúne 68.
Como o mapa desenha um marcador por estabelecimento, parte dos pontos fica sobreposta, e a concentração visual é menor que a real.

% >>> SUBSTITUIR pelos resultados reais da conferência da amostra de cerca de 30 endereços (seção de testes).
Na conferência da amostra de 30 endereços, o deslocamento mediano entre o ponto atribuído e o endereço informado foi de [XX] metros, e [XX] dos 30 pontos ficaram a menos de [XX] metros do endereço.

\subsection{Avaliação heurística}

% >>> SUBSTITUIR este texto genérico e a Tabela 3 pelos resultados reais da inspeção heurística.
A inspeção foi feita por três avaliadores, entre eles o orientador Paulo Henrique Rodrigues (\href{mailto:496762@pucminas.edu.br}{496762@pucminas.edu.br}) e dois colegas com experiência em desenvolvimento de \textit{software}.
Na consolidação das inspeções, foram identificados [N] problemas distintos, sendo [N] leves, [N] médios e [N] graves.
A Tabela \ref{tab:heuristica} apresenta os problemas apontados por mais de um avaliador, com a heurística de \citeonline{nielsen1994} relacionada e a gravidade atribuída.

\begin{table}[htbp]
    \centering
    \caption{Problemas de usabilidade apontados por mais de um avaliador}
    \label{tab:heuristica}
    \small
    \begin{tabular}{p{6.2cm}p{3.4cm}cc}
        \hline
        \textbf{Problema encontrado} & \textbf{Heurística} & \textbf{Gravidade} & \textbf{Avaliadores} \\
        \hline
        [Descrição do problema 1] & [H\# -- nome da heurística] & [Grave] & [3] \\
        \relax[Descrição do problema 2] & [H\# -- nome da heurística] & [Médio] & [2] \\
        \relax[Descrição do problema 3] & [H\# -- nome da heurística] & [Leve] & [2] \\
        \hline
    \end{tabular}
    \par\vspace{3pt}\textbf{\footnotesize Fonte: Elaborado pelo autor, a partir do registro dos avaliadores.}
\end{table}

As heurísticas com mais ocorrências foram [heurística A] e [heurística B], concentradas principalmente na tela [Inventário/Mapa].
Os avaliadores também destacaram como pontos positivos [pontos positivos da interface].
Os problemas graves foram [descrição resumida], e a correção sugerida pelos avaliadores foi [sugestão de correção].

\subsection{Discussão}

Os quatro objetivos específicos foram atendidos, com a ressalva de que a coordenada vem apenas do CEP, sem junção por endereço completo.

Em relação aos trabalhos relacionados, a taxa de 82,1\% de estabelecimentos com coordenada é menor que a de Presidente Prudente (97,4\%), onde \citeonline{silva2024} geocodificaram o endereço padronizado em um \textit{software} de geoprocessamento, e aqui a coordenada é a de um endereço do mesmo CEP.
Diferentemente do Mapa de Empresas \cite{brasil2026}, que entrega totais agregados, a aplicação lista e localiza cada estabelecimento.

Do ponto de vista prático, a aplicação permite a cidadãos, profissionais e empresários ver onde as empresas de TI se concentram em Contagem e consultar cada uma delas, e, como as coordenadas são aproximadas, a interface informa a data da fotografia e avisa que a posição é aproximada.
